\documentclass[12pt]{article}
\usepackage{graphicx} % Required for inserting images
\usepackage[margin=0.5in]{geometry}
\usepackage{amsmath, amssymb}
\usepackage{mathrsfs}


\usepackage{soul}


% Dr. David Brown Commands
\newcommand{\ds}{\displaystyle}
\newcommand{\R}{\mathbb{R}}
\newcommand{\N}{\mathbb{N}}
\newcommand{\Q}{\mathbb{Q}}
\newcommand{\C}{\mathbb{C}}


% Commands to get script r
\usepackage{calligra}
\DeclareMathAlphabet{\mathcalligra}{T1}{calligra}{m}{n}
\DeclareFontShape{T1}{calligra}{m}{n}{<->s*[2.2]callig15}{}
\newcommand{\scripty}[1]{\ensuremath{\mathcalligra{#1}}}

% Unit commands
\newcommand{\degree}{\text{°}}
\newcommand{\unitSpeed}{\frac{\text{m}}{\text{s}}}
\newcommand{\unitAcceleration}{\frac{\text{m}}{\text{s}^2}}

% Constant commands
\newcommand{\avogadro}{6.022\times10^{23}}

\newcommand{\boltzmannConstK}{1.381\times10^{-23}\frac{\text{J}}{\text{K}}}
\newcommand{\boltzmannConstInEV}{8.617\times10^{-5}\frac{\text{eV}}{\text{K}}}
\newcommand{\planckConst}{6.626\times10^{-34}\text{J}\cdot\text{s}}

\newcommand{\atmP}{1.01325\times10^5\,\text{Pa}}

% Commands to easily write partial derivatives
\newcommand{\partDeriv}[2]{\frac{\partial #1}{\partial #2}}
\newcommand{\doublePartDeriv}[2]{\frac{\partial^2 #1}{\partial^2 #2}}


\title{Problem 6.16}
\author{Gabriel Decker}


\begin{document}



\maketitle
\begin{flushleft}

Starting off, we know the following definition of $\overline{E}$ and $Z$.

\begin{equation}
    \overline{E}=\frac{1}{Z}\sum_sE(s)e^{-\beta E(s)}
\end{equation}
\begin{equation}
    Z=\sum_se^{-\beta E(s)}
\end{equation}
where $\beta=1/kT$.\\~\\

We'll start with the $Z$ equation and take the derivative $\partDeriv{}{\beta}$.
$$\partDeriv{Z}{\beta}=\partDeriv{}{\beta}\left(\sum_se^{-\beta E(s)}\right)$$
This derivative passes through to each individual term in the summation. Each one will bring its respective $E(s)$ up front, along with a negative sign. This yields a new summation:
$$\implies\partDeriv{Z}{\beta}=-\sum_sE(s)e^{-\beta E(s)}$$
This looks like the above definition for $\overline{E}$, but it's missing a few elements.
$$\implies-\frac{1}{Z}\partDeriv{Z}{\beta}=\frac{1}{Z}\sum_sE(s)e^{-\beta E(s)}$$
Now, this is the definition of $\overline{E}$. Therefore
$$\implies\overline{E}=-\frac{1}{Z}\partDeriv{Z}{\beta}$$
This expression can be simplified by remembering the derivative of $\ln(x)$ and the chain theorem.
$$\implies\overline{E}=-\partDeriv{}{\beta}\ln(Z)$$



\end{flushleft}

\end{document}
